\chapter{Research Methodology}
\label{ch:methodology}

\epigraph{The artifact and the theory it embodies are evaluated together.}{after
\citet{marchsmith1995}}

\section{Design Science Research}
\label{sec:m-dsr}

This thesis is a work of \emph{design science}: it produces and evaluates an
artifact in order to learn something general, rather than observing a
pre-existing phenomenon. The paradigm is set out by \citet{marchsmith1995}, who
distinguish \emph{build} and \emph{evaluate} as the core design activities, and
by \citet{hevner2004}, whose three-cycle view---a \emph{relevance} cycle to the
application environment, a \emph{rigour} cycle to the existing knowledge base, and
a central \emph{design} cycle of building and evaluating---structures the work
reported here. The design-science research methodology (DSRM) of
\citet{peffers2007} supplies the process spine: problem identification,
objectives, design and development, demonstration, evaluation, and communication.
\citet{wieringa2014} frames the alternation of \emph{design problems} and
\emph{knowledge questions} that this thesis interleaves.

\Cref{tab:dsrm} maps the DSRM activities onto the chapters, so that the
methodology is legible as a thread through the document rather than a detached
preamble.

\begin{table}[t]
\centering
\caption[DSRM activities mapped to chapters]{The design-science research
methodology of \citet{peffers2007} mapped onto this thesis.}
\label{tab:dsrm}
\small
\begin{tabularx}{\textwidth}{@{}l Y l@{}}
\toprule
\textbf{DSRM activity} & \textbf{Realisation in this thesis} & \textbf{Chapter} \\
\midrule
Problem identification & Narrowness of machine language; reach without accountability; fragmentation & \cref{ch:introduction} \\
Objectives of a solution & The unified-substrate and phase-transition hypotheses; RQ1--RQ5 & \cref{ch:introduction,ch:framework} \\
Design \& development & The \lspmax\ artifact and the fusion architecture & \cref{ch:artifact,ch:fusion} \\
Demonstration & Conformance across six language domains & \cref{ch:all-language} \\
Evaluation & RQs revisited; bounded-status grading; threats to validity & \cref{ch:discussion} \\
Communication & The thesis itself, and the artifact's public surfaces & all \\
\bottomrule
\end{tabularx}
\end{table}

\section{The Artifact and Its Study}
\label{sec:m-artifact}

The central artifact is \lspmax, a law-state runtime studied as it exists in the
repository at CalVer~26.6.18 \citep{lspmax}. The thesis does not report the
construction of a greenfield system; it analyses an extant, non-trivial codebase
as the instantiation of \cref{def:lawstate}, and reasons from its design. The
study proceeded by (i)~reading the artifact's constitution and architecture
documents, (ii)~tracing its protocol surface, conformance types, and enforcement
mechanisms to source with file-and-line precision, and (iii)~relating each
mechanism to the conceptual framework of \cref{ch:framework}. Where the artifact
references sibling systems that were not present in the study environment---the
\code{wasm4pm} execution engine and its type-authority crate among them---those
systems are characterised from the artifact's own interfaces and documentation,
and the gap is recorded as \unknownstat\ rather than papered over.

\section{Evaluation Strategy}
\label{sec:m-eval}

Design-science artifacts admit several evaluation methods \citep{hevner2004};
this thesis uses two. The first is \emph{analytical}: the artifact's mechanisms
are examined against the properties \cref{ch:framework} requires of a law-state
runtime---three-valued verdicts, read-only observation, receipt-bound admission,
object-centric logging---and judged present, partial, or absent. The second is
\emph{illustrative scenarios}: \cref{ch:all-language} works each of six language
domains through the pipeline of \cref{fig:pipeline} to show, by construction, that
the event-log abstraction and conformance surface apply. Illustrative scenarios
demonstrate feasibility and expose limits; they do not establish field efficacy,
and \cref{ch:discussion} is explicit that broad empirical validation across
deployed systems remains \openstat.

\section{Bounded-Status Epistemics}
\label{sec:m-epistemics}

The artifact under study enforces an epistemic discipline through its
constitution and its \code{anti-llm-cheat-lsp} canary: claims must carry
\emph{bounded status}, ``victory language'' is forbidden, and an assertion of
admission requires a \emph{receipt}---a content-addressed record binding a digest,
a digest algorithm, a boundary, a checkpoint, the raw command, and, where
required, a negative-control result \citep{lspmax}. This thesis adopts the same
discipline reflexively, for two reasons. Methodologically, it is the honest
register for design-science claims, which are warranted to varying and stateable
degrees. Substantively, it would be incoherent to describe a conformance artifact
in the unbounded language that artifact refuses.

Concretely, every substantive finding is graded with one of the bounded statuses
in \cref{tab:status}, and the three conformance verdicts never collapse
(\cref{prin:three-valued}). Numerical claims about the future are marked
\forecast\ and attributed to their forecasting body
(\cref{ch:vision}); quotations from living specifications are pinned to the
revision in force.

\begin{table}[t]
\centering
\caption[Bounded-status vocabulary]{The bounded-status vocabulary used to grade
findings, adopted from the artifact's constitution \citep{lspmax}.}
\label{tab:status}
\small
\begin{tabularx}{\textwidth}{@{}l Y@{}}
\toprule
\textbf{Status} & \textbf{Meaning in this thesis} \\
\midrule
\admitted   & Evidence is present and has been checked against the relevant criterion. \\
\refused    & Evidence is present and is contrary to the claim. \\
\unknownstat & Admissibility cannot yet be determined; a precondition or trace is missing. \\
\candidate  & A design or claim is plausible and partially supported, not yet admitted. \\
\partialstat & Admitted in part; a stated remainder is outstanding. \\
\blocked    & Determination is prevented by an external precondition. \\
\openstat   & A known gap that no current evidence closes. \\
\forecast   & A projection about the future; not a measurement. \\
\bottomrule
\end{tabularx}
\end{table}

\section{Threats to Validity (Preview)}
\label{sec:m-threats}

The threats are stated in full in \cref{sec:disc-threats}; they are previewed
here so that the reader carries them through the design chapters. \emph{Construct
validity} is at risk where the water-to-steam metaphor is taken too literally;
\cref{sec:fw-phase} confines it to an order-of-magnitude, same-character claim.
\emph{Internal validity} is at risk where the artifact's documented status is
taken for its runtime behaviour; the study relies on source and specification,
not on a live deployment exercised across all domains, and says so.
\emph{External validity} is at risk in generalising from one artifact and six
illustrative domains to ``all language''; the event-log criterion
(\cref{def:eventlog}) is the explicit boundary of the claim. \emph{Reliability}
is supported by tracing claims to primary sources and to file-and-line locations,
so that a later reader can re-examine them.
